\documentclass[12pt,a4paper]{article}

\input{core/setup} 

\begin{document}

\input{core/header}


\section*{Atividade 1}
\textit{Considerando o contexto definido na Prática 3, Quais seriam as perguntas do tipo “por que”, "para quê" que poderiam ser utilizadas nesta situação para identificar mais detalhes sobre a solução a ser desenvolvida? Devem ser definidas no mínimo 20 perguntas.}

\begin{enumerate}
    \item Quais pontos positivos e negativos da atual prática ou sistematização da gestão de pessoas?
    \item Quais indivíduos ficam responsáveis pela gestão de pessoas?
    \item Qual parte do sistema atual é automatizada e qual parte é manual?
    \item Quais tipo de atividades são necessárias agendamentos e qual a dificuldade da descentralização dessas agendas e as pessoas responsáveis por cada agenda?
    \item Qual o impacto da descentralização de informação e informalidade no registro das demandas atuais da empresa? Cite situações na qual a descentralização causou transtornos?
    \item Qual o maior "gargalo" ou dificuldade quando você precisa localizar uma cláusula específica em um contrato antigo? 
    \item Qual o ciclo de vida de um contrato, quais pessoas são responsáveis pela ratificação e como é feito o envio a cada parte importante do processo?
    \item Qual a maneira que a chefia faz a verificação da produtividade de cada funcionário, existem registros para criar métricas de desempenho?
    \item Quais processos administrativos atualmente são realizados de maneira manual e quanto tempo gasta na realização deles?
    \item Como é feita a busca de um funcionário específico? Existe divisão por setor de atuação? Quem tem acesso a essa informação? E se um funcionário mudar de setor?
    \item Quais critérios objetivos definem que uma tarefa deve ir para a Pessoa A e não para a Pessoa B? 
    \item Como o sistema deveria lidar com folgas, horas extras e férias, seria interessante trazer registros dessas informações?
    \item Faz sentido existir um meio de comunicação pelo sistema entre os seus usuários para além da descrição de uma demanda?
    \item Faz sentido atribuição de prazos de atividade e status de concluído não concluído atrasado suspenso
    \item Cada pessoa no sistema deve ter um cargo específico de modo que certos cargos possam receber certas categorias de demanda, essas demandas estão restritas a pessoas de um setor ou de vários setores de modo que uma demanda possa ser alocada a um setor responsável ou pessoa responsável
    \item Qual o nível de conhecimento dos funcionários na empresa têm com tecnologias de gestão administrativa (Excel, Word, Kanban)?
    \item Como são realizados avisos gerais internos (reuniões, treinamentos, cursos) atualmente?
    \item Como seria uma interface que se aproxima das rotinas de trabalho dos usuários do sistema?
    \item Por que o atraso na geração de relatórios impacta a tomada de decisão da diretoria?
    \item Poderia citar os 3 erros mais comuns que ocorrem no preenchimento manual atual?
    \item Por que a produtividade da equipe não é visível com as ferramentas atuais?
    \item Para que o histórico de produtividade individual será consultado?

\end{enumerate}

\section*{Atividade 2}
\textit{Considerando a Atividade 1 da Prática 4 e a Prática 3, selecione e execute um dos métodos de Elicitação de Requisitos vistos em sala, com o objetivo de:}

\begin{itemize}[before=\itshape]
    \item Identificar itens que descrevem requisitos funcionais
    \item Identificar itens que descrevem escolhas de projeto (não funcionais)
\end{itemize}

\textbf{Método escolhido: Entrevista/Conversa}

{\small\textit{(Transcrito por TurboScribe.)}}

\begin{entrevista}
    \textbf{[Cleyton]} \\
    Não existia sistema. Aí você me pergunta como é que ela acompanhava. Planilha da Excel. Ela cadastrava todos os contratos com a linha da Excel, fazia as colunas lá, com data responsável pelo contrato, a descrição do contrato fiscal, o número do processo do contrato, as observações, tudo ela colocava lá na planilha da Excel, em várias planilhas, porque as planilhas eram divididas por áreas, ela não podia pegar contrato de manutenção elétrica com contrato de fornecimento de café, por exemplo, não tem contato para isso, nem com contrato de internet, por exemplo.

    \textbf{[Felipe]} \\
    Então os contratos eram divididos por assuntos?

    \textbf{[Cleyton]} \\
    Sim, por áreas.

    \textbf{[Felipe]} \\
    Por áreas, né? Planilha da Excel. Quem é que ficava responsável por fazer a gestão dos contratos, centralizar em uma só pessoa ou em várias pessoas?

    \textbf{[Cleyton]} \\
    Aí você tem que entender um pouco como é que um contrato funciona numa instituição pública.

    \textbf{[Felipe]} \\
    Como é que um contrato funciona numa instituição pública?

    \textbf{[Cleyton]} \\
    Existe uma licitação, existe um edital, tem uma comissão que gera uma solicitação para o edital de licitação ser publicado em um diário oficial e as empresas se interessam e participam com um pregão eletrônico para demonstrar interesse. E aquela que cumprir com todos os requisitos do contrato, do edital e tiver o menor preço ou sempre o menor preço, ela ganha com a licitação, aí ela se compromete a fornecer o serviço ou fornecer o produto conforme pagamento, entendeu? Mas aí o contrato também tem que ter o fiscal e o gestor do contrato. O gestor é o que capacita, que levanta as características, as necessidades e o fiscal é o que efetivamente acompanha o cumprimento ou a execução do que foi licitado.

    \textbf{[Felipe]} \\
    Entendi, então tem alguém que vai fazer um acompanhamento de cada parte do contrato, fazendo análise do contrato, tem alguém que vai fazer a distribuição das etapas do contrato, né?

    \textbf{[Alessandro]} \\
    Sim. É isso? Então no caso, o contrato de início, ele só tem um gestor e com o decorrer dele vai ter um fiscal só?

    \textbf{[Cleyton]} \\
    O fiscal. Também o fiscal, um gestor e um fiscal. E posteriormente também um suplente do fiscal, tem um fiscal adjunto.

    \textbf{[Felipe]} \\
    Mas é só uma pessoa que é a responsável, tipo, é um gestor para um contrato, um fiscal para um contrato e um suplente para um contrato. Tipo, ele não pode os três pertencem ao mesmo contrato?

    \textbf{[Cleyton]} \\
    Não. É, exato. Os três pertencem ao mesmo contrato.

    \textbf{[Felipe]} \\
    Entendi. Só que a PGE possui dezenas, até centenas de contratos, né?

    \textbf{[Cleyton]} \\
    Sim.

    \textbf{[Felipe]} \\
    Então fica vários contratos para um gestor, vários contratos para um... a legislação contratual impede que um mesmo gestor assuma muitos contratos, né?

    \textbf{[Cleyton]} \\
    Tem um limite aí. Porque se ele assumir muitos contratos, Ela tem que verificar aí quantos... uns dez contratos são o máximo.

    \textbf{[Felipe]} \\
    Ah, tá.

    \textbf{[Cleyton]} \\
    Tem que ser certo.

    \textbf{[Felipe]} \\
    Como é que funcionava antes, quando era feito o plano de Excel, sobre a organização dos contratos? Era, tipo, se perdia muito contrato, se perdia muito prazo?

    \textbf{[Cleyton]} \\
    Isso acontecia frequentemente, porque o prazo era acompanhado no calendário, né? Como não havia nenhum tipo de automação em relação ao mesmo de prescrição ou validação do contrato, qualidade do prazo contratual, ela tinha que abrir a planilha e manualmente verificar se estava com isso no prazo.

    \textbf{[Alessandro]} \\
    Eu ia dar aquela ideia de falar contigo, que sabe...

    \textbf{[Cleyton]} \\
    Sim, qual é o teu nome?

    \textbf{[Felipe]} \\
    Felipe.

    \textbf{[Cleyton]} \\
    Felipe, quem poderia responder com exatidão essa pergunta de como é que era o setor de contratos antes da entrada do sistema é a própria Bruna, que é a gerente do contrato. Eu não sei te dizer perfeitamente, eu imagino que era... Muitas planilhas, alguma coisa no Word, anotações e post-it, né? Esses negocinhos aqui, colocar no computador pra avisar o prazo e ela olhar o calendário todos os dias e quando ela anotar assim, ah, vem esse contrato aqui, aqui e aqui. Esse era o controle que eu imagino que ela tivesse antes.

    \textbf{[Felipe]} \\
    Então a Bruna, ela que era a gerente... A gerente do contrato. Então a gente pode discutir o paradigma depois da pós-existência do sistema de contratos, do SIGESCON. O que mudou quando se adicionou o SIGESCON em relação a automações? Existem processos que antes eram manuais que foram automatizados?

    \textbf{[Cleyton]} \\
    Com certeza. Principalmente na parte de cadastro, acompanhamento, edição e alerta. Que é muito importante pra não perder o prazo de renovação de contrato.

    \textbf{[Felipe]} \\
    Como é que foi feita a solução sobre a questão dos contratos por área? A centralização dos contratos, tipo, eles foram organizados no banco de dados? Eles foram estruturados?

    \textbf{[Cleyton]} \\
    Foram organizados através do banco de dados e tabelas. E o ideal é você abrir lá o diagrama de tirada de relacionamento das tabelas do banco de dados, no próprio post-it, eu imagino, e abrir lá o PGAdmin e olhar lá como é que estão organizadas as tabelas.

    \textbf{[Felipe]} \\
    Então foi estruturado de forma, basicamente, foi estruturada. As informações que antes eram soltas em tabelas. Exato. Outra questão também, quais são, tipo, as principais soluções do SIGESCON, o que ele exatamente faz?

    \textbf{[Cleyton]} \\
    Boa pergunta, cara. Você tem que abrir porque, boa pergunta, você tem que abrir o SIGESCON e olhar lá. A versão não tem as funcionalidades, fazer o readme.

    \textbf{[Felipe]} \\
    Isso é uma coisa interessante. Tem documentação? Atualmente, quais são os pontos, tipo, que não estão fechados em relação ao SIGESCON e ao sistema de contratos?

    \textbf{[Cleyton]} \\
    Vários, vários.

    \textbf{[Felipe]} \\
    O que que falta? Tipo, que se tem demanda constante, que se tem solicitações, que... Tem algum to-do aí para o que falta fazer?

    \textbf{[Cleyton]} \\
    Nosso sistema está em constante evolução. Não, isso aí é instalação da infraestrutura. Tecnologia de visado, funcionalidade. Talvez você me perguntou o que que ele faz. Ele tem gestão de contratos, gerenciamento de identidade, controle de acesso, interface expansiva. Tem do BEC também.

    \textbf{[Felipe]} \\
    Mas ele, no geral, ele dá só um visor. Geral, né?

    \textbf{[Cleyton]} \\
    Geral. Da funcionalidade do sistema em si, e não as características do que que ele tem que fazer. Isso aí é das interfaces. Volta lá. Ali no menu, no SIGESCON. No sistema mesmo. Lá na... No sistema. No site. Lá no site. Ali, ó. Lá do esquerdo tem o dashboard de contratos, modalidade de contratos, os contratados, usuários. Ah, ele faz gestão de usuários também, né? É, Gestão de usuários.

    \textbf{[Felipe]} \\
    Como é que os usuários, eles responderam a um ambiente tipo... A um ambiente de usuário, tipo, ao SIGESCON? Essa mudança. O que ele chama de usuário aí, que na verdade são os atores dos contratos, né?

    \textbf{[Cleyton]} \\
    Isso. São os fiscais. Mas não são eles que usam. Quem faz o acompanhamento contratual, o fiscal, por exemplo, não necessariamente ele tem que acessar o sistema para preencher o cadastro. Isso pode ser feito pelo PAE. Por exemplo, quando tem uma nota fiscal para ser atestada, quando tem algum relatório, ele é enviado pelo processo administrativo. Então, tipo, aí é mais para a gente... Ele concentra os dados cadastrágeos também, dados contratuais.

    \textbf{[Felipe]} \\
    Então, o SIGESCON também é que ele é mais para consulta, enquanto o cadastro é feito pela gerência.

    \textbf{[Alessandro]} \\
    Mas não é só a gerência.

    \textbf{[Cleyton]} \\
    É só a gerência que tem permissão para cadastrar contratos, por exemplo. Enquanto você cadastra o contrato, tem algumas dúzias de lacunas a serem preenchidas lá.

    \textbf{[Alessandro]} \\
    No caso, os prazos de contratos até então, eles estão sendo preservados?

    \textbf{[Cleyton]} \\
    Sim, tem sido preservado, com certeza. Até porque muitos deles são vitais para o funcionamento da turbuladoria.

    \textbf{[Felipe]} \\
    O SIGESCON, ele é muito utilizado no sistema PGE? Ele é um sistema muito utilizado?

    \textbf{[Cleyton]} \\
    Diariamente, pelo setor competente. Sim, sem dúvida.

    \textbf{[Felipe]} \\
    E os usuários, eles lidam bem com a interface dele?

    \textbf{[Cleyton]} \\
    Vou te dar uma dica em relação ao uso, que é bastante enganosa. Por exemplo, a GDAP, que é o sistema que nós desenvolvemos, passamos aí mesmo a desenvolver. Ela é um sistema pouco usado. Está bom de visão? Ele é usado somente uma semana. Só que o impacto que ele causa é gigantesco, porque é financeiro. Sim. A partir da avaliação, o valor que é avaliado lá, que é resultante, ele reflete durante o futuro, durante os 4 próximos meses, durante o quadrimestre inteiro. Então, a questão de ele ser muito usado ou não, isso é enganoso. Entendeu? Porque o sistema pode ser pouco usado, só que tem um impacto gigante. Pensa na urna eletrônica. Tem um sistema lá que é só usado no dia da eleição, mas antes dele, né? A cada 4 anos, a cada 2 anos. Mas o impacto dele é gigantesco. Ele tem que ser testado, retestado, testado novamente, questão de segurança e validação.

    \textbf{[Felipe]} \\
    Então, é muito enganoso. Então, no caso do SIGESCON, ele seria tipo a união das duas coisas, né? Porque ele é tanto usado frequentemente quanto tem um impacto financeiro grande. Porque o setor de contratos está perdendo os maiores setores, né?

    \textbf{[Cleyton]} \\
    Dinheiro não é tudo, cara. Por exemplo, o LEX, que é o Sistema Informativo de Leis e Decretos e Portarias, ele não traz retorno financeiro nenhum, não diretamente. É um sistema quase que bibliográfico, né? Você consulta as leis nele. Então, não é só o dinheiro que nos move, também a questão de facilidade de acesso a informações. É para ajudar no trabalho do dia a dia dos setores.

    \textbf{[Felipe]} \\
    E no caso, quem tem o acesso aos contratos, ele é um acesso público? Tipo, as pessoas em geral podem ver os contratos?

    \textbf{[Cleyton]} \\
    O próprio site da PGE tem um portal de transparência pública, onde aparecem listados, a tabela de todos os contratos vigentes. Até por questão de tribunal de contas, do Estado, do TCE, ele anualmente vem aqui à PGE realizar a prestação de contas. Isso tudo bate, pede anotação fiscal, RDC, atestos.

    \textbf{[Felipe]} \\
    O sistema, ele tem acesso externo?

    \textbf{[Cleyton]} \\
    O SIGESCON, ainda não. Só o portal da transparência.

    \textbf{[Felipe]} \\
    Normalmente, quem mexe com o SIGESCON, ele tem um conhecimento... Qual o nível de conhecimento dele em tecnologia? De quem usa o SIGESCON?

    \textbf{[Cleyton]} \\
    Nível médio e nível superior. Nível de usuário? 

    \textbf{[Alessandro]} \\
    No caso, a questão que ele está falando, eu acho mais no sentido de conhecimento não acadêmico, mas a adaptação dele com a internet, com a tecnologia em si.

    \textbf{[Cleyton]} \\
    Como a maioria dessas funcionalidades foram criadas por demanda, então foi a própria usuária, a Bruna, que solicitou o que ela queria que contivesse no sistema. Então, ele é \textit{tailor-made}. Ele é feito para alfaiate, para roupa caber no cliente.

    \textbf{[Felipe]} \\
    Mas ele é voltado para os contabilistas, que vão trabalhar com essa questão?

    \textbf{[Cleyton]} \\
    Sim, para os administradores. Vão trabalhar também para o setor jurídico, ele também foi feito moldado ao setor jurídico. Sim, ele tem uma ferramenta de relatórios, de todo tipo e jeito, justamente para auditoria e conferência.

    \textbf{[Felipe]} \\
    E no nível médio, seria por conta de licitações de tipo contrato, contratar uma pessoa para fazer um serviço para o Estado?

    \textbf{[Alessandro]} \\
    É, para essa ideia. Ele também gera algum tipo de análise de desempenho em si. Mas, por exemplo, o que ele cadastra são usuários. Então, assim, devemos ver a quantidade de contratos de um usuário. Ele vê, assim, fulaninho tem menos contrato, então ele está tendo um desempenho menor.

    \textbf{[Cleyton]} \\
    Essa questão da contabilidade, de quantos contratos estão associados a um gestor, a um fiscal, é importante para, justamente, associações futuras. O Humberto está sobrecarregado de gestão de contratos, então ele não pode mais pegar nenhum contrato como gestor. Ele repassa para uma outra pessoa. E como fiscal vice-versa.

    \textbf{[Felipe]} \\
    Essa barreira de negócios é importante, né?

    \textbf{[Cleyton]} \\
    É importante. Por questão jurídica, porque não pode ocupar. Tem que consultar na lei como ele funciona. Na tela de legislação contratual, ela está sob constante mudança. Teve uma recente, aí, de uns anos atrás, de como é que a lei de contratos e contratações e licitações funciona. A lei das licitações. Isso muda o tempo inteiro. Todo ano tem mudança. Em relação ao valor, em relação ao prazo, em relação ao tipo de emprego. Porque quando a licitação é muito grande, é feito a convite, é carta convite, não é aberto ao público, entendeu? Por exemplo, a construção de uma hidrelétrica. Você não pode abrir ao público. Você já tem que ir direcionado para os empreiteiros, para quem pode participar da licitação. Ela tem que comprovar a capacidade de cumprimento da licitação. Então, já tem que ter um histórico.

    \textbf{[Felipe]} \\
    Como é feita a comunicação? É via edital? Nessa questão da automação, por exemplo. É criado o contrato aqui? O contrato, ele já tem um modelo automático? Ele tem que ser gerado? Tem que ser escrito?

    \textbf{[Cleyton]} \\
    Tem vários modelos de contrato. Eu uso até dizer algumas dezenas de modelos de contrato, dependente do tipo de produto e serviço.

    \textbf{[Felipe]} \\
    Eles são centralizados no sistema ou eles são...

    \textbf{[Cleyton]} \\
    Essa parte de edição do edital de licitação, ele não está no sistema. Isso aí é feito manualmente pela uma comissão. A gente tem que abrir um Word lá e editar todos os parágrafos e requisitos para ser licitado. Licitação é algo complicado. Demora meses até anos. E pode ser embargado. A licitação está no meio, está para terminar, chega uma empresa lá, entra com o recurso e diz que ela foi injustiçada e ela não concorda com aquela decisão. Aí paralisa a licitação até o juiz julgar e decidir se prossegue ou não com aquela licitação.

    \textbf{[Felipe]} \\
    Existe o histórico dos contratos, por exemplo? Porque você falou que existe uma \textit{timeline} do processo do contrato. O contrato é escrito, ele é pedido para ser validado...

    \textbf{[Cleyton]} \\
    Sim, ele é assinado em cartório e tudo mais.

    \textbf{[Felipe]} \\
    O sistema atualmente tem uma gestão desse histórico do processo?

    \textbf{[Cleyton]} \\
    Boa pergunta, cara. Você tem que abrir o sistema e pedir e olhar. Por isso que eu lhe digo, a continuação dessa conversa aqui tem que ser realizada com a Bruna, porque ela é a operadora máster do sistema. Ela pode abrir lá e mostrar detalhe por detalhe o que ela pediu, o que quer cumprir, e justamente perguntar para ela o que ela quer saber, o que falta melhorar no sistema e que novas funcionalidades ela deseja. Coisa que eu não parei para sentar e conversar porque ele está soberbado de outras tarefas.

    \textbf{[Alessandro]} \\
    Está chegando na parte de funcionalidade que a gente não fala quase, então... Obrigado aí, Cleyto. De verdade, ajudou demais.
\end{entrevista}

{\small\textit{(Transcrito por TurboScribe.)}}

\subsubsection*{Requisitos Funcionais (RF)}
    \begin{enumerate}
        \item \textbf{Cadastro de Contratos:} O sistema deve permitir o registro completo de contratos, incluindo data, responsável, descrição do fiscal, número do processo e observações.
        \item \textbf{Categorização por Áreas:} O sistema deve organizar os contratos por eixos temáticos (Manutenção, Fornecimento, Internet) para evitar a mistura de dados.
        \item \textbf{Gestão de Atores do Contrato:} O sistema deve permitir a associação de três perfis distintos a cada contrato como um Gestor, um Fiscal e um Suplente/Adjunto.
        \item \textbf{Sistema de Alertas de Prazos:} O sistema deve emitir notificações automáticas para evitar a perda de prazos de renovação e vencimento contratual.
        \item \textbf{Gestão de Identidade e Acesso:} O sistema deve controlar permissões, garantindo que apenas a Gerência realize o cadastro de novos contratos.
        \item \textbf{Relatórios de Auditoria:} O sistema deve gerar relatórios diversificados para conferência jurídica, contábil e para prestação de contas ao TCE.
        \item \textbf{Dashboard Estatístico:} O sistema deve prover um painel visual para acompanhamento de modalidades, contratados e quantitativo de usuários.
        \item \textbf{Validação de Regra de Negócio (Carga de Trabalho):} O sistema deve monitorar a quantidade de contratos por gestor para evitar sobrecarga (limite sugerido de 10 contratos).
    \end{enumerate}
\subsubsection*{Requisitos Não Funcionais (RNF)}
    \begin{enumerate}
        \item \textbf{Banco de Dados Relacional:} O sistema deve utilizar uma estrutura de tabelas relacionadas, preferencialmente em PostgreSQL (conforme menção ao PGAdmin e diagramas de relacionamento).
        \item \textbf{Interface Responsiva:} O sistema deve possuir uma "interface expansiva" que se adapte às necessidades de visualização do usuário.
        \item \textbf{Restrição de Acesso Externo:} O acesso administrativo deve ser restrito ao ambiente interno da Procuradoria, sem exposição para a internet.
        \item \textbf{Modelo \textit{Tailor-Made} (Sob Medida):} O desenvolvimento deve ser focado nas demandas específicas coletadas com a operadora máster (Bruna), priorizando a usabilidade para os gestores.
    \end{enumerate}

\section*{Atividade 3}
\textit{Reescreva os seguintes requisitos para remover ambiguidades:}

\subsection*{\textit{O sistema bibliotecário deve ser fácil de usar.}}
    O sistema bibliotecário deve permitir que um usuário realize a reserva de um livro em no máximo 4 cliques a partir da tela inicial e que usuários novos concluam o cadastro básico em menos de 5 minutos sem auxílio manual.

\subsection*{\textit{O sistema deve fornecer um serviço confiável a todas as classes de usuários.}}
    O sistema deve garantir uma disponibilidade (\textit{uptime}) de 99,9\% durante o horário comercial (08:00 às 18:00) e possuir um tempo médio de recuperação de falhas críticas inferior a 30 minutos.

\subsection*{\textit{O sistema deve permitir uma resposta rápida a todos os pedidos de informação.}}
    O sistema deve processar consultas ao banco de dados e retornar os resultados na interface do usuário em um tempo de latência não superior a 500 milisegundos para 95\% das requisições simultâneas.

\end{document}